% ===== 原文 pp.55--58 =====

% [原文 p.55]
$Z_{\mathrm{inst}}$ 的 $k$-瞬子贡献是如下矩阵模型围道积分
\begin{equation}
  \resizebox{0.985\textwidth}{!}{$\displaystyle Z^{(k)}_{\mathrm{inst}} = \int \dd\phi\,
  \dfrac{\prod_{i,F,I}(m^i_F-\phi^i_I)\ \prod_{i=1}^{p-1}\left[\prod_{I,J}S(\phi^{i+1}_J-\phi^i_I)\prod_{A,J}(\phi^{i+1}_J-a^i_A+\epsilon_1+\epsilon_2)\prod_{I,B}(a^{i+1}_B-\phi^i_I)\right]}
  {\prod_i\left[\left(\frac{\epsilon_1\epsilon_2}{\epsilon_1+\epsilon_2}\right)^{k_i}\prod_{I\neq J}S(\phi^i_I-\phi^i_J)\prod_{A,I}(\phi^i_I-a^i_A+\epsilon_1+\epsilon_2)(a^i_A-\phi^i_I)\right]}$}
  \tag{5.33}
\end{equation}
其中 $S(\phi)=-(\phi+\epsilon_1)(\phi+\epsilon_2)/[\phi(-\phi-\epsilon_1-\epsilon_2)]$，各指标遍历由其所出现的位置自然决定的范围（$1\le i\le p$，$1\le F\le M_i$，$1\le I\le k_i$，$1\le J\le k_{i+1}$，$1\le A\le N_i$，$1\le B\le N_{i+1}$），而 $\dd\phi$ 表示所有 $\dd\phi^i_I$ 的乘积。分子中的第一个因子刻画基本超多重态的效应，分子的其余部分来自双基本表示超多重态，分母则来自矢量多重态。

被积函数可以从图 3 所示的 IIA 弦论构造来理解。矢量多重态的贡献（分母）来自同一区间内不同 D0 膜之间伸展的弦（$1/S$ 因子）以及 D0 膜与 D4 膜之间伸展的弦（其余因子）。基本超多重态的贡献来自连接 D0 膜与 D6 膜的弦。跨越一个 NS5 膜两侧、连接 D0 膜或 D4 膜的弦给出双基本表示超多重态的贡献，其中略去了 D4--D4 相互作用，因为它们已在 $Z_{\mathrm{one-loop}}$ 中计入。

% [Tutor Note: p.55 (5.33) 逐因子核对：分母里 $(\frac{\epsilon_1\epsilon_2}{\epsilon_1+\epsilon_2})^{k_i}$ 是 $k_i$ 个 D0 膜的整体因子（注意 $\epsilon_1+\epsilon_2$ 在分母，其量纲为质量幂次 $k_i$，与分子中 $\phi$ 的幂次相消）。核对 $S$ 的归一化：原文分母写法 $\phi(-\phi-\epsilon_1-\epsilon_2)$ 里两个负号相消，故 $S(\phi)=\frac{(\phi+\epsilon_1)(\phi+\epsilon_2)}{\phi(\phi+\epsilon_1+\epsilon_2)}$；当 $\phi\to\infty$ 时 $S\to1$，这正是紫外归一化（所以 (5.33) 的测度在 $\phi$ 很大时趋于自由矩阵模型），而当 $\phi\to0$ 时 $S\sim\epsilon_1\epsilon_2/\phi$ 发散，对应 (5.34) 的极点簇。另注意正文所说的“$1/S$ 因子”正对应该分母：矢量多重态贡献含 $1/S$。分子第一因子中的指标 $(i,F,I)$ 实际只在 $M_i\ne0$ 的 $i$ 上非空；原文省略了这一点（即“indices run over the ranges that are natural given where they appear”）。]

同样的公式也可以用数学上严格的 ADHM 构造从第一性原理得到。(5.33) 中的围道使得它包围如下极点
\begin{equation}
  \{\phi^i_I\mid 1\le I\le k_i\} = \{a^i_A+(r-1)\epsilon_1+(s-1)\epsilon_2\mid (r,s)\in\lambda^i_A\}
  \tag{5.34}
\end{equation}
其中 $\lambda^i_A$（$1\le i\le p$，$1\le A\le N_i$）遍历所有这样的杨图集合：其格子总数等于瞬子数，即 $\sum_{A=1}^{N_i}|\lambda^i_A|=k_i$。极点集合 (5.34) 来自一个称为 Jeffrey--Kirwan（JK）留数规定的处方，它由等变局域化推出 [210]。该处方也可以从前述 ADHM 超对称矩阵模型的局域化得到 [241--244]。

\begin{exercise}{5.4}
把这些公式特殊化到具有 $N_f=4$ 个味的 U(2) SQCD（$p=1$，$N_1=2$，$M_1=4$）。注意，通过平移 Coulomb 分支参数 $a^1_A$（$A=1,2$），四个质量参数 $m^1_B$（$B=1,\dots,4$）都可以被平移。这是如下事实的一个影子：在 U(2) 理论中我们只有 SU($N_f$) 味对称性，而不像 SU(2) SQCD 那样有 SO($2N_f$)。计算 $Z_{\mathrm{inst}}$ 中的单瞬子贡献。并开始计算双瞬子贡献以体会其复杂度：写出被积函数并列出极点。
\end{exercise}

% [Tutor Note: p.55 关于 Exercise 5.4 的提示（非原文）：$p=1$ 时 (5.33) 中 $\prod_{i=1}^{p-1}$ 为空（等于 1），只剩基本超多重态因子与矢量多重态因子；单瞬子时 $k_1=1$，$I=J=1$ 故 $I\ne J$ 的 $S$ 因子也不出现，被积函数是 $1/\prod$ 形态的有理函数，极点由 (5.34) 给出，即 $\phi=a_A$（$r=s=1$，空杨图）、$\phi=a_A+\epsilon_1$、$\phi=a_A+\epsilon_2$ 等单格杨图位置。]

\para{U(1) 因子。} 该理论的膜实现需要规范化 U(1) 对称性，以便把 SU($N_i$) 补成 U($N_i$) 规范群。在此过程中，U(1) 质量参数
% [原文 p.56]
成为动力学矢量多重态标量。幸好，$Z_{\mathrm{inst}}$ 是在这些标量取常数值时计算的，因此我们可以简单地把这些 Coulomb 分支参数取为我们出发时的 U(1) 质量参数。

\begin{figure}[htbp]
\centering
\includegraphics[width=0.60\textwidth]{fig3.png}
\figtitle{3}{$U(4)\times U(4)\times U(3)$ 规范理论的膜构造。最左侧的 U(4) 因子带两个基本超多重态（由横向 D6 膜插入，图中记作 $\otimes$），最右侧的 U(3) 因子带一个基本超多重态。在这个例子中，每个规范群都是渐近自由的。}
\end{figure}

规范化带来的主要麻烦效应是：从 U(1) 规范因子会引入多余的瞬子贡献。在与 2d CFT 共形块对照时发现不匹配、困惑了一周之后，[5] 的作者们提议把这个多余的 U(1) 瞬子贡献除掉，示意性地\footnote{欢迎提供参考文献：最好能把这一手续在瞬子配分函数一侧写清楚，我很可能漏掉了正确的文献。}
\begin{equation}
  Z_{\mathrm{inst}}\Big[\prod \mathrm{SU}(N_i)\Big] = Z_{\mathrm{inst}}\Big[\prod \mathrm{U}(N_i)\Big]\Big/\left(\mathrm{U}(1)\ \text{factor}\right).
  \tag{5.35}
\end{equation}
这个 U(1) 因子在 CFT 一侧更容易被单独挑出。SU($N$) 规范群的线性箭图的瞬子配分函数被期望与 $W_N$ 代数（$N=2$ 时为 Virasoro）的共形块相配，而 U($N$) 规范群则对应 $W_N$ 乘以 Heisenberg 手征代数的共形块。于是这两种情形的差别（更确切地说是比值）由 Heisenberg 手征代数的共形块给出，即手征顶点算子的自由场关联函数。这一观点在 $N=2$ 时已由 [14] 理解，后来在 [245] 中被推广。对这里所考虑的线性箭图，该因子示意性地为
\begin{equation}
  \left(\mathrm{U}(1)\ \text{factor}\right) \simeq \prod_{1\le i\le j\le p}(1-z_i\cdots z_j)^{c_{ij}}
  \tag{5.36}
\end{equation}
其中某些指数 $c_{ij}$ 是质量参数的二次型。

% [Tutor Note: p.56 (5.36) 的结构核对：指数 $c_{ij}$ “对质量参数是二次型”意味着 $\log$ 后的 U(1) 因子为 $\sum_{i\le j}c_{ij}\log(1-z_i\cdots z_j)$，而 Heisenberg 自由场的手征顶点算子关联函数型如 $\langle\prod_a \ee^{\ii q_a X}(z_a)\rangle=\prod_{a<b}(z_a-z_b)^{q_aq_b}$，其指数确实是荷的双线性，故 $c_{ij}$ 与质量（即 U(1) 荷）成二次关系是自洽的。另一可核对锚点：$p=1$ 时 (5.36) 只剩 $(1-z_1)^{c_{11}}$，正是把 U(1) 因子按 $\log(1-z)$ 展开的形态。]

\para{重整化方案：一个例子。} 鉴于提取 U(1) 因子的问题，以及关于哪些理论可以用膜构造处理的额外限制，

% [原文 p.57]
人们寻求了其他方法来确定 $Z_{\mathrm{inst}}$。我们介绍其中一种方法 [96]，它例示了比较不同结果时的一个微妙之处。

带有定向折叠（orientifold）面的膜构型使人们能得到由 USp 群与 Spin 群交替构成的线性箭图的 $Z_{\mathrm{inst}}$。我们将在 7.2 小节末尾回到这类箭图及其 class S 构造。事实上，最简单的例子之一就是 SU(2) SQCD（$N_f=4$）的另一种描述：
\begin{equation}
  \vcenter{\hbox{\includegraphics[width=0.51\textwidth]{eq537.png}}}
  \tag{5.37}
\end{equation}
更一般地，USp(2) $=$ SU(2) 与 Spin(4) $=$ SU(2) $\times$ SU(2) 规范群的线性箭图等价于广义 SU(2) 箭图。在足够简单的情形如 (5.37) 中，瞬子配分函数也可以借助 U(2) ADHM 构造并除掉 U(1) 因子 (5.35) 来计算。令人惊讶的是，两种方法给出的、按指数化规范耦合 $q_{\mathrm{Lag}}=\ee^{2\pi\ii\tau_{\mathrm{Lag}}}$ 的幂展开的级数并不相同，这乍看之下意味着矛盾！

同样的问题很久以前就在比较来自 4d $\cN=2$ 理论不同构造的 SW 解时出现过。它被 [246] 通过指出耦合定义中的重整化方案含混而解决。我们在比较 SU(2) SQCD 拉格朗日量中出现的规范耦合 $\tau_{\mathrm{Lag}}$ 与该理论 class S 构造中出现的位形参数 $q=\ee^{2\pi\ii\tau}$ 时，已在 (4.20) 中遇到过这种含混，并发现了一个关系 $\tau=\tau_{\mathrm{Lag}}+\cdots$，其中含一个常数平移以及无穷多个瞬子修正\footnote{在 AGT 的语境中，一个较早指出这些瞬子修正的文献是 [247]，另可参见 [248] 的一个 F 理论推导。}。

让我们集中于（质量形变的）SCFT，例如 SU(2) $N_f=4$ SQCD，它们的规范耦合不跑动，并假设只有一个规范耦合以简化表达式。在两个领头阶相一致的重整化方案中，规范耦合常数 $\tau,\tilde\tau$ 之间的关系为
\begin{equation}
  \tilde\tau = \tau+\sum_{n\ge0}c_n\ee^{2\pi\ii n\tau},
  \tag{5.38}
\end{equation}
其中与 $\tau$ 无关的系数 $c_n$ 是无量纲的。在具体情形中这些系数与 Coulomb 分支及质量参数无关；若假设两个方案在整个参数空间上都在领头阶一致，一个可能的论证是：系数 $c_n$ 应当是 Coulomb 分支与质量参数的有界函数，从而由解析函数的 Liouville 定理必为常数。

U(2) 方法与 USp(2)--Spin(4) 方法对 $Z_{\mathrm{inst}}$ 的计算正是以这种方式相关联的，这一点在 [96] 中已得到理解。考虑具有 $N_f=4$ 个无质量超多重态的 SU(2) SQCD。箭图描述 (5.37) 对应于把 SQCD 视为 $\cX(\fso(4))=\cX(\fsu(2))\otimes\cX(\fsu(2))$ 带上 $\Z_2$ 自同构扭曲缺陷（twist defect）的约化，或者视为 $\cX(\fsu(2))$ 的约化：
\begin{equation}
  \vcenter{\hbox{\includegraphics[width=0.62\textwidth]{eq539.png}}}
  \tag{5.39}
\end{equation}

% [Tutor Note: p.57 (5.37) 与 (5.39) 的原文都是示意图，按 STYLE_GUIDE §6bis 采用裁图（eq537.png、eq539.png）。读法：(5.37) 左边是“方框--椭圆--方框”= Spin(4)--USp(2)--Spin(4) 的线性箭图（方框为 Spin 节点，椭圆为 USp 节点），右边把每个 Spin(4)=SU(2)×SU(2) 拆成上下两个 SU(2) 方框，中间用交叉（X）表示相邻节点之间的双基本连线；这正是“广义 SU(2) 箭图”的图例约定。(5.39) 中 $\mathrm{T}[\fso(4),\text{带切割线的圆柱}]=\mathrm{T}[\fsu(2),\text{四点球面}]$ 表示 tinkertoy：so(4) 曲线（圆柱，两穿刺点加分支切割）的双覆盖正是 su(2) 曲线（四点球面），与 p.58 开头一句完全对应。]

% [原文 p.58]
$\fso(4)$ 曲线的双覆盖（带一条分支切割的圆柱）是 $\fsu(2)$ 曲线（四点球面）。被切割圆柱的模参数 $q_{\mathrm{so}(4)}$ 与红外耦合之间的关系，由与四点球面的模参数 $q_{\mathrm{su}(2)}$ 相同的考量确定（见 (4.20)）。由此可以把两者相互关联起来，
\begin{align}
  q_{\mathrm{so}(4)} &= \sqrt{16\lambda(2\tau_{\mathrm{IR}})} = 16\ee^{\pi\ii\tau_{\mathrm{IR}}}+O(\ee^{3\pi\ii\tau_{\mathrm{IR}}}), \nonumber\\
  q_{\mathrm{su}(2)} &= \lambda(\tau_{\mathrm{IR}}) = q_{\mathrm{so}(4)}\Big(1+\frac{q_{\mathrm{so}(4)}}{4}\Big)^{-2} = 16\ee^{\pi\ii\tau_{\mathrm{IR}}}-128\ee^{2\pi\ii\tau_{\mathrm{IR}}}+O(\ee^{3\pi\ii\tau_{\mathrm{IR}}}).
  \tag{5.40}
\end{align}
在比较 Nekrasov 瞬子配分函数时，必须考虑计数参数 $q_{\mathrm{su}(2)}$ 与 $q_{\mathrm{so}(4)}$ 之间的这个非平凡映射。

% [Tutor Note: p.58 (5.40) 的数值锚点（可自行展开验证）：取 $\lambda(\tau)=16q-128q^2+704q^3-\cdots$（$q=\ee^{\pi\ii\tau}$），则 $\lambda(2\tau_{\mathrm{IR}})=16\ee^{2\pi\ii\tau_{\mathrm{IR}}}-128\ee^{4\pi\ii\tau_{\mathrm{IR}}}+\cdots$，于是 $\sqrt{16\lambda(2\tau_{\mathrm{IR}})}=16\ee^{\pi\ii\tau_{\mathrm{IR}}}(1-8\ee^{2\pi\ii\tau_{\mathrm{IR}}}+\cdots)^{1/2}=16\ee^{\pi\ii\tau_{\mathrm{IR}}}-64\ee^{3\pi\ii\tau_{\mathrm{IR}}}+\cdots$，与原文的其余项阶数 $O(\ee^{3\pi\ii\tau_{\mathrm{IR}}})$ 一致。再用 $(1+x)^{-2}=1-2x+3x^2-\cdots$（$x=q_{\mathrm{so}(4)}/4=4\ee^{\pi\ii\tau_{\mathrm{IR}}}+\cdots$）：$q_{\mathrm{so}(4)}(1-8\ee^{\pi\ii\tau_{\mathrm{IR}}}+\cdots)=16\ee^{\pi\ii\tau_{\mathrm{IR}}}-128\ee^{2\pi\ii\tau_{\mathrm{IR}}}+\cdots$，正是原文第二式的 $-128$ 系数------这是判断该式没有笔误的最好办法。]

\subsection{用局域化切割}

\para{没有拉格朗日量时的超对称局域化。} 到目前为止，我们展示的超对称局域化都是针对拉格朗日理论的，因为该技术依赖于路径积分表述。在没有路径积分表述时，我们无法像以前那样把作用量形变为 $S\to S+tQV$，但仍然可以把期望值 $\langle\cO\rangle$ 形变为 $\langle\cO\ee^{-tQV}\rangle$。与 (5.12) 中相同，只要超荷 $Q$ 在所关心的态上不自发破缺、且 $Q\cO=0$，这个期望值就与 $t$ 无关。在 $t\to+\infty$ 极限下，可观测值 $\langle\cO\rangle$ 应当在某种意义上约化为“$(QV)_{\text{bosonic}}$ 的零点”，只要这种说法能被定义。

在下文中\footnote{这些想法大多已见于文献，但可能没有以这种方式被收集过。}，我们只用矢量多重态的形变项，就把一般 class S 理论的配分函数分解为对 $C$ 的任意裤子分解的简单构件之积。另一个无关的方面是：用一个限制在 $S^3_b\subset S^4_b$ 上的形变项，我们把配分函数因子化为 $q$ 的（反）全纯函数之积的积分，正如 (5.28) 那样，这对任何 3d 限制为拉格朗日量的理论都成立。

\para{通过局域化矢量多重态来切割 $C$。} 高阶 class S 理论通常是非拉格朗日量的，其获得方式是规范化某些孤立 SCFT（称为 tinkertoy，最著名的是 $T_N$）的公共味对称性。然而，规范化对称性是用标准矢量多重态来完成的，而这些多重态由标准路径积分描述，其规范耦合 $\tau$ 与拉格朗日情形相同。这就启示我们把超对称局域化应用于这些矢量多重态。

让我们更明确地说出“规范化某理论 T（tinkertoy 之积）的味对称群 $G$”是什么意思。首先把 T 中的 $G$ 对称性流多重态耦合到一个非动力学矢量多重态 $(A,\lambda,\phi)$。配分函数 $Z_{\mathrm{T}}[A,\lambda,\phi]$ 依赖于这个背景矢量多重态，而其各分量可以取任意值（常数或非常数）。规范化就是对矢量多重态场作路径积分，于是规范化后理论的配分函数为
\begin{equation}
  Z_{\mathrm{T,gauged}} = \int [\mathrm{D}A\,\mathrm{D}\lambda\,\mathrm{D}\phi]\,\ee^{-S_{\mathrm{SYM}}[A,\lambda,\phi]}Z_{\mathrm{T}}[A,\lambda,\phi],
  \tag{5.41}
\end{equation}
